\documentclass{article}

\textwidth=6.5in
\textheight=8.5in
\voffset=-54pt
\oddsidemargin=0in
\evensidemargin=0in
\setlength{\parskip}{8pt}
\setlength{\parindent}{0pt}

\usepackage{workbook}

\usepackage[sols]{handout}

\begin{document}

\begin{center}
\large\textsc{Problem Set 18 - Relative Growth Rates}
  
      \pspace{12pt}
  \begin{problemonly}\large Name: \rule{5cm}{0.15mm}\\ \end{problemonly}
\end{center}


\begin{grading}
\begin{enumerate}
    \item 12 points
    \item 10 points
    \item 10 points
    \item 10 points
    \item 15 points
    \item 15 points
\end{enumerate}
Total: 72 points.
\end{grading}

\begin{framed}
   \textit{Throughout this PSET (and unit), you will work to evaluate different limits. When performing these evaluations, you may need to use L'H\^opital's Rule, or you may be able determine the limit by using your knowledge of relative growth rates. Assessing which method a given limit may require is an important part of the problem-solving process!} 
\end{framed}

\begin{enumerate}

\item 

\begin{grading}
12 points. 3 points per part. They can use limits to justify their answer here, or they can appeal to relative growth rates.
\end{grading}

% Meredith / Meghan?

  Recall that, if $f(x)$ and $g(x)$ are functions that are both positive
  when $x$ is large, we say that $g(x)$ \textbf{grows more quickly than}
  $f(x)$ if $\displaystyle \lim_{x \to \infty} \frac{f(x)}{g(x)} = 0$ (or,
  equivalently, if $\displaystyle \lim_{x \to \infty} \frac{g(x)}{f(x)} =
  {\infty}$).
  For each pair of functions below, determine which grows more quickly.
  Justify each of your answers using a limit.

  For example: $f(x)=x\ln x$ grows faster than $g(x) = x$, since
  $\displaystyle \lim_{x\to\infty} \frac{x\ln x}{x} = \displaystyle
  \lim_{x\to\infty} \ln x = \infty$.

  \begin{enumerate}
  \item
\begin{problem}[\vfill]
      $\displaystyle f(x) = x\ln x$ or $\displaystyle g(x) = x^2$
    \end{problem}

    \begin{solution}
      Let's calculate $\displaystyle \lim_{x \to \infty}
      \frac{f(x)}{g(x)}$ (it doesn't matter whether we calculate this
      limit or $\displaystyle \lim_{x \to \infty} \frac{g(x)}{f(x)}$; both
      will give us the information we want).  
      \begin{align*}
        \lim_{x \to \infty} \frac{x \ln x}{x^2} & \nolheq \lim_{x \to
          \infty} \frac{\ln x}{x} \text{ (just canceling an $x$ from the
          numerator and denominator)} \\
        &\lheq \lim_{x \to \infty} \frac{\ \frac{1}{x}\ }{1} \\
        & \nolheq 0,
      \end{align*}
      so $\boxed{x \ln x \ll x^2}$.
    \end{solution}

  % \item
  %   \begin{problem}[\vfill]
  %     $f(x) = 10^{1000} x \ln x$ or $\displaystyle g(x) = x^2$.  (Does
  %     your answer change if you change $10^{1000}$ to a larger constant?) 
  %   \end{problem}

  %   \begin{solution}
  %     The limit $\displaystyle \lim_{x \to \infty} \frac{10^{1000} x \ln
  %       x}{x^2}$ is just $10^{1000}$ times the limit we calculated in
  %     {{{{\#2}}(a)}}; that limit was 0, so $\displaystyle
  %     \lim_{x \to \infty} \frac{10^{1000} x \ln x}{x^2} = 0$ as well.
  %     Therefore, $\boxed{10^{1000} x \ln x \ll x^2}$.  More generally, if
  %     $k$ is any constant then $kx \ln x \ll x^2$. 
  %   \end{solution}

  \item
    \begin{problem}[\vfill]
      $\displaystyle f(x) = x^2e^x$ or $\displaystyle g(x) = x^2+e^x$
    \end{problem}

    \begin{solution}
      Let's calculate $\displaystyle \lim_{x \to \infty}
      \frac{g(x)}{f(x)}$, which is an $\frac{\infty}{\infty}$ limit.
      L'H\^opital's Rule works, but it's easier to evaluate the limit
      using just algebra:
      \begin{align*}
        \lim_{x \to \infty} \frac{x^2 + e^x}{x^2 e^x} & = \lim_{x \to
          \infty} \left( \frac{x^2}{x^2 e^x} + \frac{e^x}{x^2 e^x} \right)
        \\
        &= \lim_{x \to \infty} \left( \frac{1}{e^x} + \frac{1}{x^2}
        \right) \\
        &= 0,
      \end{align*}
      so $\boxed{x^2 + e^x \ll x^2 e^x}$.       
    \end{solution}

  \item % Janet

    \begin{problem}[\vfill\newpage]
      $\displaystyle f(x) = x^{1000} e^x$ or $g(x) = e^{2 x}$
    \end{problem}

    \begin{solution}
      Since $x^{1000} \ll e^x$, you should expect $x^{1000} e^x \ll e^x
      e^x = e^{2x}$.  Let's check:
      \begin{align*}
        \lim_{x \to \infty} \frac{f(x)}{g(x)} & \nolheq \lim_{x \to
          \infty}
        \frac{x^{1000} e^x}{e^{2x}} \\
        \shortintertext{Using exponent rules, we can rewrite this as} %
        & \nolheq \lim_{x \to \infty} \frac{x^{1000}}{e^{x}} \\
        & \lheq \lim_{x \to \infty} \frac{1000 x^{999}}{e^{x}} \\
        & \lheq \lim_{x \to \infty} \frac{1000 \cdot 999 x^{998}}{e^{x}}
      \end{align*}
      If we continue using L'H\^opital's Rule, the denominator will always
      be $e^{x}$, while the numerator will eventually
      (after using L'H\^opital's Rule another 998 times) just be a number.
      This limit will be 0, so $\boxed{x^{1000} e^x \ll e^{2x}}$.
    \end{solution}

  \end{enumerate}

\item

\begin{grading}
10 points. 9 points for part (a), and 1 point for part (b). It's fine if students use the ``shortcut" of relative growth rates for limits as $x\to\infty$, but they \textit{cannot} do that for limits as $x\to -\infty$!
\end{grading}

\begin{enumerate}
  \item Find all horizontal asymptotes of each
    function. Show the work leading to your answer.

    \begin{enumerate}
    \item
      \begin{problem}[\vfill]
        $f(x) = \dfrac{e^{2x} + 4 e^x + 3}{e^{2x} + 1}$
      \end{problem}

      \begin{solution}
        As $x \to -\infty$, $e^x$ and $e^{2x}$ both approach 0, so
        $\boxed{ \lim_{x \to -\infty} f(x) = \frac{3}{1} = 3}$.

        On the other hand, to evaluate the limit as $x \to \infty$, we can
        make use of relative growth rates.  In the sum $e^{2x} + 4 e^x +
        3$, $e^{2x}$ is the fastest growing term.  In the sum $e^{2x} +
        1$, $e^{2x}$ is the faster growing term.  So, $\boxed{\lim_{x \to
            \infty} f(x) = \lim_{x \to \infty} \frac{e^{2x}}{e^{2x}} =
          1}$.         
      \end{solution}

    \item
      \begin{problem}[\vfill]
        $g(x) = \dfrac{(2x)^4 + 7x^2 + 4x + 8}{16 x^4 - 2x^2 + 2}$
      \end{problem}

      \begin{solution} 
        As $x \to \infty$, the fastest growing term in the sum $(2x)^4 +
        7x^2 + 4x + 8$ is $(2x)^4 = 16 x^4$, and the fastest growing term
        in the sum $16 x^4 - 2x^2 + 2$ is $16 x^4$, so $\boxed{\lim_{x \to
            \infty} g(x) = \lim_{x \to \infty} \frac{16 x^4}{16x^4} = 1}$.

        The situation is exactly the same as $x \to -\infty$, because the
        $x^4$ terms are still the fastest growing; therefore, $\boxed{
          \lim_{x \to -\infty} g(x) = \lim_{x \to -\infty} \frac{16
            x^4}{16 x^4} = 1}$.
      \end{solution}

    \item
      \begin{problem}[\vfill]
        $h(x) = \dfrac{2x^2 + 10x + 8}{x^4 + 2x^2 + 2}$
      \end{problem}

      \begin{solution} 
        Now, $\boxed{ \lim_{x \to \infty} h(x) = \lim_{x \to \infty}
        \frac{2x^2}{x^4} = 0}$ and $\boxed{\lim_{x \to -\infty} h(x) =
        \lim_{x \to -\infty} \frac{2x^2}{x^4} = 0}$. 
      \end{solution}

    \end{enumerate}

  \item
    Match the three functions above with their graphs.

    \begin{imagetable}{3}{\Alph}
      \begin{tikzpicture}
        \begin{axis}[xtick=\empty, width=2.25in, ymin=-1, ymax=6, ytick={1,
            2, ..., 5}]
          \addplot+[domain=-10:10, samples=401] {((2*x)^4+7*x^2+4*x+8)/(16*x^4-2*x^2+2)};
        \end{axis}
      \end{tikzpicture}
      &
      \begin{tikzpicture}
        \begin{axis}[xtick=\empty, width=2.25in, ymin=-1, ymax=6, ytick={1,
            2, ..., 5}]
          \addplot+[domain=-10:10] {(e^(2*x)+4*e^x+3)/(e^(2*x)+1)};
        \end{axis}
      \end{tikzpicture}
      &
      \begin{tikzpicture}
        \begin{axis}[xtick=\empty, width=2.25in, ymin=-1, ymax=6, ytick={1,
            2, ..., 5}]
          \addplot+[domain=-10:10, samples=201] {(x^2+10*x+8)/(x^4+2*x^2+2)};
        \end{axis}
      \end{tikzpicture}      
    \end{imagetable}

    \begin{solution} % Jason Bland
      Given our limits, i must match (B), ii must match (A), and iii must
      match (C).
    \end{solution}

  \end{enumerate}

    \pspace{\stretch{0.25}}
  \pspace{\newpage}

% \item 

% \begin{grading}
% 10 points. 4 for part (a), and 6 for part (b). (3 for each discontinuity)
% \end{grading}


% Limits can be useful when we are trying to sketch a function.
%   In this problem, you'll look at this idea more deeply. 
%   \begin{enumerate}
%   \item In each part below (i - iv), you are given a description of the
%     limiting behavior of a function $f(x)$ at $x = a$.  From this, what
%     can you conclude about the graph of $f$?  (Select from one of the
%     choices, A - D, below.)

%     \begin{enumerate}
%     \item
%       \begin{problem}[\vfill]
%         $\displaystyle \lim_{x \to a^-} f(x) = \infty$ or $-\infty$ and
%         $\displaystyle \lim_{x \to a^+} f(x) = \infty$ or $-\infty$
%       \end{problem}

%       \begin{solution}
%         $f$ has a vertical asymptote, choice \fbox{D}.
%       \end{solution}

%     \item

%       \begin{problem}[\vfill]
%         $\displaystyle \lim_{x \to a} f(x)$ exists but is not equal to
%         $f(a)$.
%       \end{problem}

%       \begin{solution}
%         This is what happens in choice \fbox{B}. 
%       \end{solution}

%     \item
%       \begin{problem}[\vfill]
%         $\displaystyle \lim_{x \to a} f(x)$ exists and is equal to
%         $f(a)$. 
%       \end{problem}

%       \begin{solution}
%         This is what happens in choice \fbox{A}.
%       \end{solution}

%     \item
%       \begin{problem}[\vfill]
%         $\displaystyle \lim_{x \to a^-} f(x)$ and $\displaystyle \lim_{x
%           \to a^+} f(x)$ both exist, but they are not equal to each other.
%       \end{problem}

%       \begin{solution}
%         This is what happens in choice \fbox{C}.
%       \end{solution}

%     \end{enumerate}

%     \begin{minipage}{1.0\linewidth}
%       Possible conclusions (each is illustrated with two examples):
%       \begin{multicols}{2}
%         \begin{enumerate}[itemsep=8pt,label=\Alph*.]
%         \item $f(x)$ is \textbf{continuous} at $x = a$.  
%           \begin{center}
%             \begin{tikzpicture}
%               \begin{axis}[xtick={3}, ytick=\empty, xticklabels={$a$},
%                 domain=-2:4, width=1.7in]
%                 \addplot+{(x - 1)*(x+1)*(x-3)+6};
%               \end{axis}
%             \end{tikzpicture}
%             \hspace{0.15in}
%             \begin{tikzpicture}
%               \begin{axis}[xtick={1}, ytick=\empty, xticklabels={$a$},
%                 domain=-2:4, width=1.7in, ymin=0]
%                 \addplot+{abs(x-1) + 4};
%               \end{axis}
%             \end{tikzpicture}
%           \end{center}

%         \item $f(x)$ has a \textbf{hole} at $x = a$.  (A hole is more
%           officially called a \uline{removable discontinuity}.)
%           \begin{center}
%             \begin{tikzpicture}
%               \begin{axis}[xtick={3}, ytick=\empty, xticklabels={$a$},
%                 domain=-2:4, width=1.7in]
%                 \addplot+{(x - 1)*(x+1)*(x-3)+6};
%                 \addplot[only marks, mark=*, fill=white]
%                 coordinates { (3, 6) };
%               \end{axis}
%             \end{tikzpicture}
%             \hspace{0.15in}
%             \begin{tikzpicture}
%               \begin{axis}[xtick={1}, ytick=\empty, xticklabels={$a$},
%                 domain=-2:4, width=1.7in, ymin=0]
%                 \addplot+{abs(x-1) + 4};
%                 \addplot[only marks, mark=*, fill=white]
%                 coordinates { (1, 4) }; 
%                 \addplot[only marks, mark=*, fill=black]
%                 coordinates { (1, 2) }; 
%               \end{axis}
%             \end{tikzpicture}
%           \end{center}

%         \item $f(x)$ has a \textbf{jump} at $x = a$.
%           \begin{center}
%             \begin{tikzpicture}
%               \begin{axis}[xtick={3}, ytick=\empty, xticklabels={$a$},
%                 domain=-2:4, width=1.7in, cycle list=black]
%                 \addplot+[domain=-2:3]{(x - 1)*(x+1)*(x-3)+6};
%                 \addplot+[domain=3:4] {16-(x-3)^2};
%                 \addplot[only marks, mark=*, fill=white]
%                 coordinates { (3, 6) (3, 16)};
%               \end{axis}
%             \end{tikzpicture}
%             \hspace{0.15in}
%             \begin{tikzpicture}
%               \begin{axis}[xtick={1}, ytick=\empty, xticklabels={$a$},
%                 domain=-2:4, width=1.7in, ymin=0, cycle list=black]
%                 \addplot+[domain=-2:1]{abs(x-1) + 4};
%                 \addplot+[domain=1:4] {x + 1};
%                 \addplot[only marks, mark=*, fill=white]
%                 coordinates { (1, 4) }; 
%                 \addplot[only marks, mark=*, fill=black]
%                 coordinates { (1, 2) }; 
%               \end{axis}
%             \end{tikzpicture}
%           \end{center}

%         \item $f(x)$ has a \textbf{vertical asymptote} at $x = a$.
%           \begin{center}
%             \begin{tikzpicture}
%               \begin{axis}[xtick=\empty, extra x ticks={1}, ytick=\empty,
%                 extra x tick labels={$a$}, domain=-2:4, width=1.7in,
%                 ymin=-2, ymax=5]
%                 \addplot+{ sin(deg(x))/(x-1)+1};
%               \end{axis}
%             \end{tikzpicture}
%             \hspace{0.15in}
%             \begin{tikzpicture}
%               \begin{axis}[xtick=\empty, extra x ticks={3}, ytick=\empty,
%                 extra x tick labels={$a$},
%                 domain=-2:4, width=1.7in, ymin=0, ymax=20,
%                 restrict y to domain=0:20]
%                 \addplot+{cos(2*deg(x))/(x-3)^2 + 1};
%               \end{axis}
%             \end{tikzpicture}
%           \end{center}          
%         \end{enumerate}

%       \end{multicols}      
%     \end{minipage}

%   \pspace{\newpage}
  
%   \item
%     \begin{problem}[\stretch{3}]
%       Let $f(x) = \dfrac{\sin (x^2 - 4)}{x^2 - 4}$.  For what values of
%       $x$ is $f(x)$ NOT defined?  At each such value, does $f$ have a
%       removable discontinuity, a jump, or a vertical asymptote?  Use
%       limits to support your answers (write down each limit you're
%       calculating, how you calculate it, and what you conclude from it).
%     \end{problem}

%     \begin{solution}
%       The only thing that would cause $f$ to be undefined is the
%       denominator being 0, which happens when $x = \pm 2$.

%       Let's look first at $x = 2$.  To decide what type of discontinuity
%       occurs there, we should calculate $\displaystyle \lim_{x \to 2}
%       f(x)$, which is a $\frac{0}{0}$ limit:
%       \begin{align*}
%         \lim_{x \to 2} \frac{\sin(x^2 - 4)}{x^2 - 4} & \lheq
%         \lim_{x \to 2} \frac{\cos(x^2 - 4) \cdot 2x}{2x} \\
%         & \nolheq \lim_{x \to 2} \cos(x^2 - 4) \\
%         & \nolheq 1
%       \end{align*}
%       Since $\displaystyle \lim_{x \to 2} f(x)$ exists, $f$ has a
%       removable discontinuity (or hole) at $x = 2$.

%       Notice that $f$ is an even function (because $f(-x) = f(x)$), so $f$
%       must also have a hole at $x = -2$.

%       So, $f$ has \fbox{removable discontinuities at $x = \pm 2$}.
%     \end{solution}

%   \end{enumerate}

\item 

\begin{grading}
10 points. 3 for part (a), and 7 for part (b), as follows: 1/1/4/1. Part (b)v. does not require any student work and should not be graded; it's just there to make sure students check their answers.
\end{grading}

\begin{enumerate}
  % \item 
  %   \begin{problem}[\vfill]
  %     Explain, in your own words, why $0 \cdot \infty$ is an indeterminate
  %     form.

  %     (If you're struggling with this, read from Example F.5 on page 1115
  %     to the end of Example F.10 on page 1117 in Gottlieb, and then
  %     try this again.)
  %   \end{problem}

  %   \begin{solution}
  %     If we have a limit $\displaystyle \lim_{x \to a} f(x) g(x)$ where
  %     $f(x) \to 0$ and $g(x) \to \infty$, that means that, for $x$ near
  %     $a$, $f(x)$ is very close to 0 and $g(x)$ is huge.  But when we
  %     multiply a tiny number (close to 0) by a huge number, the result is
  %     just somewhere in between; there's not a specific value it must be
  %     close to.
  %   \end{solution}    

  % \item The strategy for dealing with a limit $\displaystyle \lim_{x \to
  %     a} f(x) g(x)$ of the indeterminate form $0 \cdot \infty$ is to
  %   rewrite the product $f(x) g(x)$ as a quotient: \[ f(x) g(x) =
  %   \frac{f(x)}{1 / g(x)} \text{ or } f(x) g(x) = \frac{g(x)}{1 / f(x)}\]
  %   Both equations are true; often, one will be easier to work with than
  %   the other.  If you try one and it's messy, try the other.
  %   You'll use this strategy below.

  %    In this problem, you'll look at what happens to $f(x) = x \sin
  % \left( \frac{1}{x} \right)$ as $x \to \infty$.


  % \begin{enumerate}
  % \item
  %   \begin{problem}[\vfill]
  %     What is $\displaystyle \lim_{x \to \infty} x$?  (What does this make
  %     you think that $\displaystyle \lim_{x \to \infty} x \sin \left(
  %       \tfrac{1}{x}\right)$ is?)      
  %   \end{problem}

  %   \begin{solution}
  %     $\displaystyle \lim_{x\to\infty} x = \infty$.  This might make you
  %     think that $\displaystyle \lim_{x \to \infty} x \sin \left(
  %       \tfrac{1}{x}\right) = \infty$, but that would be faulty reasoning:
  %     just knowing that one piece of a limit $\to \infty$ doesn't make the
  %     entire limit $\to \infty$. 
  %   \end{solution}

  % \item
  %   \begin{problem}[\vfill\newpage]
  %     What is $\displaystyle \lim_{x \to \infty} \sin \left(
  %     \tfrac{1}{x}\right )$?  (What does this make you think that 
  %     $\displaystyle \lim_{x \to \infty} x \sin \left( 
  %     \tfrac{1}{x}\right)$ is?)
  %   \end{problem}

    % \begin{solution}
    %   As $x \to \infty$, $\frac{1}{x} \to 0$, so $\sin\left(
    %   \tfrac{1}{x}\right) \to \boxed{0}$.  This might make you think that
    %   $\displaystyle \lim_{x \to \infty} x \sin\left( \tfrac{1}{x}\right)
    %   = 0$, but that would be faulty reasoning: just knowing that one
    %   piece of a limit $\to 0$ doesn't make the entire limit $\to 0$.
    % \end{solution}

  \item
\begin{problem}[\vfill]
      To deal with the limit $\displaystyle \lim_{x \to \infty}
      x\sin(\tfrac{1}{x})$, we can rewrite it:
      \[ \lim_{x \to \infty} x\sin(\tfrac{1}{x}) = \lim_{x \to \infty} \frac{\sin(\tfrac{1}{x})}{\tfrac{1}{x}}\]  Use L'Hôpital's Rule to evaluate the limit.  (Why is it valid to use L'Hôpital's
      Rule?) 
    \end{problem}

    \begin{solution}
      $\displaystyle \lim_{x \to \infty} \frac{\sin\left(
      \tfrac{1}{x}\right)}{ \frac{1}{x}}$ is a
      $\frac{0}{0}$ limit, so we can use
      L'Hôpital's Rule:
      \begin{align*}
        \lim_{x\to\infty} \frac{\sin\left(\tfrac{1}{x}\right)}{\frac{1}{x}}
        & \lheq \lim_{x\to\infty} \frac{-\frac{1}{x^2}\cdot
          \cos\left(\frac{1}{x}\right)}{-\frac{1}{x^2}} \\
        \intertext{We can cancel $-\frac{1}{x^2}$ in the numerator and
          denominator:} %
        & \nolheq \lim_{x \to \infty} \cos\left(\tfrac{1}{x}\right) \\
        \shortintertext{As $x \to \infty$, $\frac{1}{x} \to 0$, so this limit
          is equal to} %
        & \nolheq \cos(0) \\
        & \nolheq \boxed{1}
      \end{align*}
    \end{solution}

  \item
    \begin{problem}[\stretch{0.5}]
      Another way to evaluate the limit is to make a substitution.
      Let $u=\tfrac{1}{x}$. As $x \to \infty$, $u \to 0^+$.
      We can thus rewrite the limit as
      \[\lim_{x \to \infty} \frac{\sin\left(\frac{1}{x}\right)}
      {\frac{1}{x}} = \lim_{u \to 0^+} \frac{\sin u}{u}.\]
      Find $\displaystyle \lim_{u \to 0^+} \frac{\sin u}{u}$. (We evaluated this special limit in lesson 3.1 before we knew L'Hôpital's Rule.)
    
    \end{problem}

    \begin{solution}
        We determined in workbook section 3.1 that
        $\displaystyle \lim_{x \to 0} \frac{\sin u}{u}=\fbox{1}$.
        This is true for both one-sided limits as well. Thus,
        \[\lim_{x \to \infty} \frac{\sin\left(\frac{1}{x}\right)}
      {\frac{1}{x}} = \lim_{u \to 0^+} \frac{\sin u}{u}=1.\]
    \end{solution}

  \item
    \begin{problem}
      Use Desmos to graph $y=x\sin(\tfrac{1}{x})$ and see whether your answer agrees.
    \end{problem}

    \begin{solution}
      You can see the graph on Desmos here: 
      \url{https://www.desmos.com/calculator/dnsmjdkbc5}.
      The graph seems to agree with our answer. There is a horizontal
      asymptote at $y=1$. As $x$ gets larger and larger, the 
      $y$-coordinate becomes closer and closer to $1$.
    \end{solution}
    

  % \end{enumerate}

    \end{enumerate}

\pspace{\newpage}

\iffalse
\item 

\begin{grading}
12 points. 4 points per part.
\end{grading}

    Evaluate the following limits.  (Be careful; not
    all of them are of the form $0 \cdot \infty$.)

    \probonly{}
      \begin{enumerate}
      \item

        \begin{problem}
          $\displaystyle \lim_{x \to \infty } \sqrt[3]{x} e^x$ $\vphantom{\dfrac{1}{1}}$
        \end{problem}

        \begin{solution}
          As $x \to \infty$, $\sqrt[3]{x} \to \infty$ and $e^x \to
          \infty$; therefore, this limit is $\boxed{\infty}$ (it's not of
          an indeterminate form).
        \end{solution}

      \item

        \begin{problem}
          $\displaystyle \lim_{x \to -\infty } \sqrt[3]{x} e^x$
          $\vphantom{\dfrac{1}{1}}$ 
        \end{problem}

        \begin{solution}
          As $x \to -\infty$, $\sqrt[3]{x} \to -\infty$ and $e^x \to 0$, so this
          is a $0 \cdot \infty$ limit, and we rewrite it so that we can
          use L'H\^opital's Rule:
          \begin{align*}
            \lim_{x \to -\infty } \sqrt[3]{x} e^x & \nolheq \lim_{x \to -\infty}
            \frac{\sqrt[3]{x}}{e^{-x}} \\
            \intertext{This is a $\frac{\infty}{\infty}$ limit, so now we
              can use L'H\^opital's Rule:} %
            & \lheq \lim_{x \to -\infty} \frac{\frac{1}{3}
              x^{-2/3}}{-e^{-x}} \\
            \intertext{Now, the numerator approaches 0 while the
              denominator approaches $-\infty$, so the limit is} % 
            & \nolheq \boxed{0}
          \end{align*}
        \end{solution}

      \item
        \begin{problem}
          $\displaystyle \lim_{x \to \infty} x^2 \sin \left(
            \frac{1}{x} \right)$
        \end{problem}

        \begin{solution}
          As $x \to \infty$, $x^2 \to \infty$ and $\sin \left(
            \frac{1}{x} \right) \to \sin(0) = 0$ since $\frac{1}{x} \to
          0$.  Therefore, this is a $0 \cdot \infty$ limit, so we should
          rewrite it so that we can use L'Hôpital's Rule:
          \begin{align*}
            \lim_{x \to \infty} x^2 \sin \left( \frac{1}{x} \right) &
            \nolheq \lim_{x \to \infty} \frac{\sin \left(
                \frac{1}{x} \right)}{x^{-2}} \\
            & \lheq \lim_{x \to \infty} \frac{\cos \left( \frac{1}{x}
              \right) \cdot - \frac{1}{x^2}}{- 2 x^{-3}} \\
            & \nolheq \frac{x \cos \left( \frac{1}{x} \right)}{2} \\
            & \nolheq \boxed{\infty}
          \end{align*}
        \end{solution}

      % \item
      %   \begin{problem}
      %     $\displaystyle \lim_{x \to \infty} x^2 \cos \left(
      %       \frac{1}{x} \right)$
      %   \end{problem}

      %   \begin{solution}
      %     As $x \to \infty$, $x^2 \to \infty$ and $\cos \left( \frac{1}{x}
      %     \right) \to \cos(0) = 1$, so this limit is $\boxed{\infty}$.
      %   \end{solution}

      \end{enumerate}      
      \probonly{}
\fi
 

% \item 

% \begin{grading}
%  15 points, as follows: 6/5/4
% \end{grading}

%  %(Appendix F. 22)

%   Let $f(x) = x^2 \ln x$, defined on $(0, \infty)$.
%   \begin{enumerate}
%   \item
%     \begin{problem}[\vfill]
%       Find the exact coordinates of all local extrema, and say whether
%       each is a local minimum or local maximum.
%     \end{problem}

%     \begin{solution}
%       To answer this, let's make a sign chart for $f'$.  First,
%       $f'(x) = x^2 \cdot \frac{1}{x} + 2x \ln x = x + 2x \ln x = x (1 + 2
%       \ln x)$.  This is never undefined on $(0, \infty)$, and it's equal
%       to 0 when $x = 0$ (not in the domain) or $\ln x = -\frac{1}{2}$
%       (which happens when $x = e^{-1/2} = \frac{1}{\sqrt{e}}$).

%       Now, let's test the sign of $f'$ at a value greater than $e^{-1/2}$
%       and a value less than $e^{-1/2}$.  For a value greater, we can use
%       $1$: $f'(1) = 1 > 0$.  For a value less than $e^{-1/2}$, let's use
%       $e^{-1}$: $f' \left( \frac{1}{e} \right) = \frac{1}{e} \left[ 1 + 2
%         (-1) \right] < 0$.  So, we have the following sign chart for $f'$:
%       \begin{center}
%         \begin{tikzpicture}
%           \begin{axis}[axis y line=none, x axis line style={-}, xtick={0.606531}, xticklabels={$\frac{1}{\sqrt{e}}$}, ymin=-2, ymax=2, height=1.25in, width=3in]
%             \addplot+[domain=0.406531:0.706531, very thin]{0};
%             \draw (axis cs: 0.406531, -1.5) node[right] {sign of $f'$};
%             \draw (axis cs: 0.406531, 1) node[right] {$f$};
%             \draw (axis cs: 0.556531, -1.5) node{$-$};
%             \draw[->, xshift=2pt] (axis cs: 0.426531, 1.5) -- (axis cs: 0.596531, 0.5);
%             \draw (axis cs: 0.656531, -1.5) node{$+$};
%             \draw[->, xshift=2pt] (axis cs: 0.616531, 0.5) -- (axis cs: 0.696531, 1.5);
%             \draw[->, xshift=2pt] (axis cs: 0.596531, 0.5) -- (axis cs: 0.616531, 0.5);
%           \end{axis}
%         \end{tikzpicture}
%       \end{center}
%       Therefore, $f$ has a local minimum at $x = \frac{1}{\sqrt{e}} =
%       e^{-1/2}$.  The corresponding $y$-value is $f \left( e^{-1/2}
%       \right) = e^{-1} \ln \left( e^{-1/2} \right) = - \frac{1}{2e}$.  So,
%       \fbox{the point $\displaystyle
%         \left(\frac{1}{\sqrt{e}},-\frac{1}{2e}\right)$ is a local
%         minimum}.
%     \end{solution}

  \item
\begin{problem}[\vfill]
      Find $\displaystyle \lim_{x \to \infty} x^2 \ln x$ and
      $\displaystyle \lim_{x \to 0^+} x^2 \ln x$.

      \emph{Hint:} You should only use L'H\^opital's Rule for one of these.
      % (which one?); to do that, it will help to rewrite $f(x)$ as either
      % $\dfrac{x^2}{(\ln x)^{-1}}$ or as $\dfrac{\ln x}{x^{-2}}$.  (These
      % are the same, but which one looks easier to use L'H\^opital's Rule
      % with?) 
    \end{problem}

    \begin{solution}
      As $x \to \infty$, both $x^2$ and $\ln x$ grow without bound, so 
      $\boxed{ \lim_{x\to\infty} x^2\ln x =  \infty}$. 

      As $x \to 0^+$, $x^2 \to 0$ and $\ln x \to -\infty$, so we'll
      rewrite $\displaystyle \lim_{x \to 0^+} x^2 \ln x$ and use
      L'H\^opital's Rule: 
      \[ \lim_{x \to 0^+} x^2 \ln x = \lim_{x\to0^+} \frac{\ln x}{x^{-2}}
      \lheq \lim_{x\to0^+} \frac{1/x}{-2x^{-3}} = \lim_{x\to0^+}
      \frac{x^2}{-2} = 0.\]
    \end{solution}

  % \item
  %   \begin{problem}[\vfill\newpage]
  %     Use the previous parts to sketch a rough graph of $f(x)$.  (Your
  %     graph does not need to accurately show the concavity of $f$.)
  %   \end{problem}

  %   \begin{solution}
  %     Here's the graph:
  %     \begin{center}
  %       \begin{tikzpicture}
  %         \begin{axis}[xtick={0.606531},
  %           xticklabels={$\frac{1}{\sqrt{e}}$}, ytick={-0.18394},
  %           yticklabels={$-\frac{1}{2 e}$}, cycle list=black] 
  %           \addplot+[domain=0:2] {x^2*ln(x)};
  %           \addplot+[only marks, mark=*, fill=white] coordinates
  %           {
  %             (0, 0) 
  %           };            
  %         \end{axis}
  %       \end{tikzpicture}
  %     \end{center}
  %     \checkmark\
  %       The key features your graph should have are:
  %       \begin{itemize}[nosep]
  %       \item It decreases on $\left( 0, \frac{1}{\sqrt{e}} \right)$ and
  %         increases on $\left( \frac{1}{\sqrt{e}}, \infty \right)$.
  %       \item $\displaystyle \lim_{x \to 0^+} f(x) = 0$ and $\displaystyle
  %         \lim_{x \to \infty} f(x) = \infty$. 
  %       \end{itemize}
  %   \end{solution}

  % \end{enumerate}

\item 

\begin{grading}
5 points.
\end{grading}

  \begin{problem}
    \textbf{Reflect Back.} As $x\to\infty$, the functions
    $\ln x$, $x^n$, and $e^x$ all increase without bound, but at different rates.  In our class, we use the notation $\ll$ to describe one function growing faster than another.

    Here are some common functions listed according to their relative
    rates of growth:
    \[ 2^x \ll 3^x \ll 10^x \ll x^x \] Insert into this list the functions
    $e^x$, $x^3$, $\ln x$, $x$, and $\sqrt{x}$.

    (Having a general idea of the relative growth rates of functions such
    as polynomials, exponentials, and logarithms is a handy skill in many
    fields.)    
  \end{problem}

  \pspace{\stretch{0.5}}

  \begin{solution}
    $\ln x \ll \sqrt{x} \ll x^3 \ll 2^x \ll e^x \ll 3^x \ll 10^x \ll
    x^x$. 
  \end{solution}




\end{enumerate}

\psreflection{}

  \end{document}